% Ficha de estudio 5 del Taller virtual ATA EPE: colisiones, contacto y fricción.
\documentclass[11pt,a4paper]{article}
% Preámbulo común de las fichas de estudio del Taller virtual ATA EPE.
\usepackage[spanish,es-noshorthands,es-tabla]{babel}
\usepackage[utf8]{inputenc}
\usepackage[T1]{fontenc}
\usepackage{lmodern}
\usepackage[margin=2.2cm]{geometry}
\usepackage[table]{xcolor}
\usepackage{amsmath,amssymb,bm}
\usepackage{tabularx,array,booktabs}
\usepackage{enumitem}
\usepackage{fontawesome5}
\usepackage[most]{tcolorbox}
\usepackage[hidelinks]{hyperref}
\usepackage{listings}

\definecolor{morado}{HTML}{2A1B7E}
\definecolor{amarillo}{HTML}{FFCF2E}
\definecolor{rojo}{HTML}{C4142C}
\definecolor{verde}{HTML}{1F9D55}
\definecolor{gris}{HTML}{5D5A73}

\newcolumntype{Y}{>{\raggedright\arraybackslash}X}
\renewcommand{\arraystretch}{1.25}
\setlist{itemsep=1pt,topsep=3pt}
\numberwithin{equation}{section}
\lstset{basicstyle=\ttfamily\small,columns=fullflexible,keepspaces=true,frame=single,rulecolor=\color{gris!40},
  literate={á}{{\'a}}1 {é}{{\'e}}1 {í}{{\'i}}1 {ó}{{\'o}}1 {ú}{{\'u}}1 {ñ}{{\~n}}1 {Ω}{{$\Omega$}}1
    {θ}{{$\theta$}}1 {λ}{{$\lambda$}}1 {Δ}{{$\Delta$}}1 {μ}{{$\mu$}}1 {·}{{$\cdot$}}1 {←}{{$\leftarrow$}}1 {§}{{\S}}1 {→}{{$\to$}}1 {≠}{{$\neq$}}1 {É}{{\'E}}1 {Á}{{\'A}}1 {∨}{{$\vee$}}1 {∧}{{$\wedge$}}1 {σ}{{$\sigma$}}1 {×}{{$\times$}}1 {≈}{{$\approx$}}1}

\newcommand{\rev}[1]{\widetilde{#1}}
\newcommand{\I}{\mathbf{I}}
\newcommand{\vect}[1]{\bm{#1}}
\newcommand{\A}{\mathcal{A}}

\newtcolorbox{meta}{colback=morado,colframe=morado,coltext=white,arc=3mm}
\newtcolorbox{lectura}{colback=morado!5,colframe=morado,fonttitle=\bfseries,title={\faBook~Lecturas}}
\newtcolorbox{ejercicio}[1]{colback=amarillo!10,colframe=amarillo!80!black,fonttitle=\bfseries,title={\faPencil*~Ejercicio #1}}
\newtcolorbox{control}{colback=rojo!5,colframe=rojo,fonttitle=\bfseries,title={\faQuestionCircle~Preguntas de control}}
\newtcolorbox{codigo}[1][De la ecuación al código]{colback=verde!6,colframe=verde,fonttitle=\bfseries,title={\faCode~#1}}
\newtcolorbox{cierre}{colback=verde!6,colframe=verde,fonttitle=\bfseries,title={\faFlagCheckered~Cierre del módulo}}
\newtcolorbox{nota}[1]{colback=amarillo!12,colframe=amarillo!80!black,fonttitle=\bfseries,title={\faLightbulb~#1}}

% \cabecera{número}{título}{subtítulo}
\newcommand{\cabecera}[3]{%
  \title{\textcolor{morado}{\textbf{Ficha #1. #2}}\\[3pt]\large Taller virtual ATA EPE · #3}
  \author{}\date{}}

\cabecera{5}{Colisiones, contacto y fricción}{Módulo 5: contacto}

\begin{document}
\maketitle
\vspace{-2.5em}

\begin{meta}
\textbf{\textcolor{amarillo}{Meta del módulo.}} Detectar cuándo y dónde se tocan dos piezas, y resolver ese contacto
con las leyes de Signorini y de Coulomb: que no se atraviesen, que no se peguen, que deslicen o se adhieran según la
fricción, y que se puedan apilar sin temblar. Es el módulo más largo. Cara pública: \textbf{la mesa} — piezas que
caen, chocan, ruedan y se apilan.
\end{meta}

\begin{lectura}
\begin{enumerate}[leftmargin=*]
  \item \textbf{Ericson, \emph{Real-Time Collision Detection}}, caps.~4 (volúmenes envolventes), 5 (pruebas entre
        primitivas, incluido SAT) y 9 (GJK).
  \item \textbf{Catto, «Iterative dynamics with temporal coherence» (GDC 2005)}: contactos persistentes, fricción
        acumulada y apilamiento estable.
  \item \textbf{Stewart y Trinkle (1996)}: el contacto con fricción como problema de complementariedad, con rigor.
  \item van den Bergen, \emph{Collision Detection in Interactive 3D Environments}, cap.~4: GJK y EPA en detalle.
  \item Baraff y Witkin (SIGGRAPH 2001), notas de contacto; Johnson, \emph{Contact Mechanics}, cap.~1 (para entender
        qué se está idealizando).
\end{enumerate}
\end{lectura}

% =================================================================================================
\section{Formas de colisión}

La malla visual (con agujeros, roscas y biseles) es demasiado compleja para chocar. Cada pieza tiene una forma
simple: unión de cajas, cilindros y cápsulas, todas convexas. Lo que define a una forma convexa es su
\textbf{función soporte}: el punto más lejano en una dirección $\vect d$,
\begin{equation}
  s_K(\vect d) = \arg\max_{\vect x\in K}\ \vect d\cdot\vect x .
\end{equation}
Para un cilindro de radio $r$, altura $h$ y eje $\hat{\vect a}$:
$s(\vect d) = \vect c + \operatorname{sgn}(\vect d\cdot\hat{\vect a})\tfrac h2\hat{\vect a} + r\,
\dfrac{\vect d_\perp}{\lVert\vect d_\perp\rVert}$, con $\vect d_\perp$ la parte de $\vect d$ perpendicular al eje.

\begin{ejercicio}{5.1}
Escribe la función soporte de una caja orientada y de una cápsula. Para una viga del kit, ¿qué error se comete al
ignorar los agujeros en la colisión? ¿Importa para el kit real?
\end{ejercicio}

% =================================================================================================
\section{Detección en dos etapas}

\subsection{Etapa gruesa}
Cada forma se encierra en una caja alineada a los ejes (AABB). Solo los pares cuyas AABB se tocan pasan a la etapa
fina. Ordenar los extremos de las cajas en un eje (\emph{barrido y poda}) evita comparar todos contra todos.
Además se descartan los pares unidos en el grafo de ensamble (ficha 4).

\subsection{Etapa fina}
\begin{itemize}[leftmargin=*]
  \item \textbf{Teorema del eje separador (SAT)}: dos convexos no se tocan si y solo si existe un eje en el que sus
        proyecciones no se superponen. Para dos cajas bastan 15 ejes: 3 normales de cada caja y 9 productos cruz de
        aristas. El eje de menor superposición da la normal y la profundidad.
  \item \textbf{GJK}: para convexos generales, trabaja con la diferencia de Minkowski $A - B$ (los convexos se tocan
        si y solo si contiene el origen) usando solo funciones soporte, $s_{A-B}(\vect d) = s_A(\vect d) - s_B(-\vect d)$.
        Construye un símplice que se acerca al origen.
  \item \textbf{EPA}: si GJK encuentra el origen dentro, expande el símplice hasta la cara más cercana al origen:
        normal y profundidad de penetración.
\end{itemize}

\begin{ejercicio}{5.2}
Demuestra que dos convexos se intersecan si y solo si $\vect 0\in A - B$. Dibuja en 2D una iteración de GJK con un
triángulo y un cuadrado.
\end{ejercicio}

\subsection{Puntos de contacto}
Una cara apoyada en otra no se toca en un punto sino en un polígono. Se recorta una cara contra la otra
(Sutherland--Hodgman) y se guardan hasta 4 puntos. Los contactos se reconocen de un paso al siguiente (\emph{contactos
persistentes}) para reutilizar sus impulsos.

\begin{nota}{En PGA}
La distancia con signo de un punto $P$ a un plano normalizado $p$ sale de $p\wedge P = d\,\I$; el plano de una cara
es un vector; recortar un polígono por un plano es evaluar ese signo en cada vértice e intersecar las aristas que lo
cambian con $L\wedge p$. Todo el recorte se escribe con las operaciones de la ficha 0.
\end{nota}

% =================================================================================================
\section{La física del contacto}

\subsection{Signorini}
En cada punto de contacto, con separación $g$ y fuerza normal $\lambda_n$:
\begin{equation}
  g \ge 0,\qquad \lambda_n \ge 0,\qquad g\,\lambda_n = 0 .
\end{equation}
O no se tocan, o se tocan y empujan; nunca tiran. Es una condición de \emph{complementariedad}.

\subsection{Coulomb}
La fuerza tangencial $\vect\lambda_t$ está en el cono de fricción y, si hay deslizamiento, se opone a él con la máxima
disipación:
\begin{equation}
  \lVert\vect\lambda_t\rVert\le\mu\lambda_n,\qquad \vect v_t \ne 0\ \Rightarrow\ \vect\lambda_t = -\mu\lambda_n\frac{\vect v_t}{\lVert\vect v_t\rVert}.
\end{equation}
Juntas, (3.1) y (3.2) forman un problema de complementariedad no lineal (Stewart y Trinkle lo plantean como LCP
aproximando el cono con una pirámide).

\subsection{Discretización con impulsos}
Se agregan a los impulsos secuenciales de la ficha 4 dos tipos de fila por contacto:
\begin{itemize}[leftmargin=*]
  \item \textbf{Normal}: impulso acumulado recortado a $P_n\ge0$; con un sesgo $-\tfrac{\beta}{h}\max(0, \delta - \delta_0)$
        para corregir la penetración $\delta$ (dejando un margen $\delta_0$ que evita el temblor), y el rebote
        $v_n^+\ge -e\,v_n^-$ si el choque es rápido.
  \item \textbf{Tangencial} (dos direcciones): impulso acumulado proyectado al disco $\lVert\vect P_t\rVert\le\mu P_n$.
\end{itemize}

\begin{ejercicio}{5.3}
Una caja que desliza con velocidad inicial $v_0$ sobre la mesa se detiene en $d = v_0^2/(2\mu g)$. Deduce por qué el
algoritmo anterior reproduce ese resultado, y escríbelo como prueba (con su tolerancia).
\end{ejercicio}

\begin{ejercicio}{5.4}
Un cubo apoyado en un plano inclinado de ángulo $\theta$: con el algoritmo, ¿para qué $\theta$ empieza a deslizar?
Compara con $\tan\theta_c = \mu$ (prueba P3).
\end{ejercicio}

% =================================================================================================
\section{Apilar sin temblar}

Lo que separa una simulación de juguete de una buena es poder apilar. Ingredientes, en orden de importancia:
contactos persistentes con arranque en caliente; margen de penetración $\delta_0$; suficientes iteraciones; y
corrección de posición separada de la velocidad (\emph{split impulse}) para que corregir la penetración no meta
energía. Los cuerpos que casi no se mueven se \emph{duermen}.

\begin{codigo}[Qué programar]
\begin{tabularx}{\linewidth}{@{}l Y@{}}
\texttt{nucleo/formas.js} & Cajas, cilindros, cápsulas y uniones de ellas; funciones soporte; AABB.\\
\texttt{nucleo/colision.js} & Barrido y poda; SAT caja–caja; GJK y EPA; contactos de cara (recorte); persistencia.\\
\texttt{nucleo/contacto.js} & Filas normal y tangencial en el solucionador de la ficha 4; rebote; margen; reposo.\\
\texttt{kit/catalogo.js} & Forma de colisión de cada pieza.\\
\texttt{app/} & La mesa: soltar, lanzar, apilar.\\
\end{tabularx}
\end{codigo}

\begin{control}
\begin{enumerate}[leftmargin=*]
  \item ¿Por qué basta conocer la función soporte de un convexo para detectar colisiones con GJK?
  \item ¿Por qué la condición de Signorini es una complementariedad y no una ecuación?
  \item ¿Qué es el principio de máxima disipación en la fricción de Coulomb?
  \item ¿Por qué se proyecta el impulso tangencial al disco $\mu P_n$ y no se resuelve exactamente?
  \item ¿Para qué sirve el margen de penetración?
  \item ¿Por qué los pares unidos en el ensamble no deben chocar?
\end{enumerate}
\end{control}

\begin{cierre}
Ejercicios 5.1--5.4 · prueba P3 y las nuevas (frenado $v_0^2/2\mu g$, rebote $e^2h$, pila de 5 vigas quieta 2\,s) ·
cara pública: la mesa con piezas que chocan y se apilan.
\end{cierre}

\end{document}
